\begin{textsample}{POS Dim 4 – 2011\_c6 – Score 2.00 – 2011\_c6/enade\_2011\_ci\_comp\_dab220f8e3.txt}  \label{ex:f4_pos_299}
CAPÍTULO 3 ANÁLISE TÉCNICA DA PROVA Este capítulo tem por objetivo apresentar o desempenho dos estudantes concluintes de Computação ( Bacharelado, Licenciatura, Sistemas de Informação e Engenharia de Computação ) no ENADE/2011.
Para isso, foram calculadas as estatísticas básicas da prova em seu todo, bem como as estatísticas dos componentes relacionadas à Formação Geral, ao de Conhecimento Específico da Área e das questões discursivas isoladamente.
Além disso, as análises das questões objetivas do componente de Conhecimento Específico são apresentadas separadamente para concluintes do Bacharelado, da Licenciatura, de Sistemas de Informação e de Engenharia de Computação, já que as questões propostas para cada grupo foram diferentes.
Nas tabelas, são apresentados o tamanho da população inscrita e de presentes, e as seguintes estatísticas das notas8 : média do desempenho na prova, erro padrão da média, desvio padrão, nota mínima, mediana e nota máxima.
As estatísticas apresentadas neste capítulo contemplam o total de estudantes concluintes da área de Computação ( Bacharelado, Licenciatura, Sistemas de Informação e Engenharia de Computação ) em 2011 do Brasil e, separadamente, por Grande Região.
Foram calculadas tendo -se em vista as seguintes agregações : ( a ) as Grandes Regiões e o país como um todo ; ( b ) a Categoria Administrativa ; e ( c ) a Organização Acadêmica.
Em relação aos gráficos de distribuição de notas, o intervalo considerado foi de 10 unidades, aberto à esquerda e fechado à direita, com exceção do primeiro intervalo, [ 0 ; 10 ], fechado em ambos os extremos.
Para os gráficos de distribuição das notas das questões discursivas, foram consideradas mais duas categorias : questão em branco e nota zero.
Todos os gráficos de distribuição de notas permitem a comparação dos resultados por habilitação : Bacharelado, Licenciatura, Sistemas de Informação e Engenharia de Computação. 3.1 ESTATÍSTICAS BÁSICAS DA PROVA 3.1.1 Estatísticas Básicas Gerais A Tabela 3.1 apresenta as estatísticas básicas da prova por grande Região dos estudantes concluintes de Computação ( Bacharelado, Licenciatura, Sistemas de Informação 8 Essas estatísticas e outras estão definidas no Capítulo 1. e Engenharia de Computação ).
A população total de inscritos foi de 27.253.
Destes, 22.819 estiveram presentes, sendo 16,3 \% o índice de não comparecimento.
A Região de maior abstenção foi a Norte ( 21,6 \% ) e a de menor abstenção foi a Sul ( 11,5 \% ).
A média das notas da prova como um todo ( nas seções seguintes serão analisados os componentes de Formação Geral e de Conhecimento Específico ) foi 31,8, sendo que os alunos das regiões Norte e Sudeste obtiveram a média mais baixa ( 31,0 ) e os das regiões Nordeste e Sul obtiveram a média mais alta ( 33,4 ).
A média na região Centro-Oeste foi 30,8.
O desvio padrão para o Brasil como um todo foi 11,6, sendo o maior desvio padrão encontrado na região Nordeste ( 12,2 ) e o menor na região Norte ( 10,4 ), indicando uma dispersão um pouco menor das notas desta última região.
A região que obteve a maior nota máxima foi a Sudeste ( 86,7 ), ao passo que a região que atingiu a menor nota máxima foi a Centro-Oeste ( 74,7 ).
A mediana do Brasil como um todo foi 31,0, sendo a maior mediana obtida na região Nordeste ( 32,7 ) e a menor obtida nas regiões Norte e Sudeste ( 30,2 ).
A nota mínima foi zero em todas as regiões, sem exceção.
O comportamento das notas dos estudantes de todo o Brasil pode ser observado no Gráfico 3.1 que apresenta a distribuição das mesmas em colunas diferentes para alunos de : Bacharelado ( azul ), Licenciatura ( verde ), Sistemas de Informação ( bege ) e Engenharia de Computação ( roxo ).
As distribuições de Engenharia de Computação e Licenciatura possuem moda no intervalo ( 30 ; 40 ].
Já para o Bacharelado e a habilitação de Sistemas de Informação o intervalo modal é ( 20 ; 30 ].
Observa -se também no Gráfico 3.1 que a distribuição das notas de Engenharia de Computação tem freqüências maiores que das demais habilitações em intervalos superiores.
O coeficiente de assimetria desta distribuição é positivo ( 0,33 ).
Apesar de haver concentração em intervalos mais baixos, a distribuição das notas dos alunos do Bacharelado tem coeficiente de assimetria também positivo ( 0,36 ), valor próximo ao de Engenharia de Computação, as duas distribuições têm cauda mais longa à direita.
As distribuições das notas de Licenciatura e Sistemas da Informação são praticamente simétricas.
Ambas têm coeficientes de assimetria próximos de zero : – 0,06 para Licenciatura e 0,19 para Sistemas da Informação.
Apenas para os cursos de Licenciatura das regiões Norte e Sul os coeficientes de assimetria são negativos, concentração maior do lado direito do histograma e mais espalhada do lado esquerdo.
Para as demais habilitações em todas as regiões a assimetria é positiva, com coeficiente variando de 0,02 a 0,46.
Os Gráficos 3.2, 3.3 e 3.4 apresentam informações referentes à média da nota final dos Participantes, desagregando os resultados de acordo com, respectivamente, as Grandes Regiões do país, a Categoria Administrativa e a Organização Acadêmica.
Os gráficos apresentam o valor da média das notas como uma barra e os extremos do intervalo de confiança de 95 \% como linhas verticais unidas por uma linha horizontal na forma da letra H maiúscula. de confiança de 95 \% como linhas verticais unidas por uma linha horizontal na forma da letra H maiúscula.
Considerando -se o gráfico de notas segundo Grande Região ( Gráfico 3.2 ), observa -se que existe diferença estatisticamente significativa ao nível de 95 \% entre a maior média, obtida nas regiões Sul e Nordeste ( 33,4 ), e as médias das outras três regiões.
Levando -se em conta os estudantes por Categorias Administrativas ( Gráfico 3.3 ), observa -se que existe diferença estatisticamente significante entre as médias das notas das IES Públicas e Privadas.
Pode ser observado ainda que a média dos alunos de IES Públicas ( 36,7 ) é maior que a dos alunos de IES Privadas ( 30,3 ).
A diferença entre IES Públicas e Privadas é maior do que as diferenças entre regiões.
Tendo como foco o Gráfico 3.4, que apresenta as notas médias das provas segundo Organização Acadêmica, constata -se que existe diferença estatisticamente significativa ao nível de 95 \% nas médias das notas dos estudantes provenientes de Universidades, Centros Universitários e Faculdades.
A maior média foi obtida pelos estudantes de Universidades ( 33,6 ), e a menor, pelos de Faculdades ( 29,4 ). 3.1.2 Estatísticas Básicas no Componente de Formação Geral A Tabela 3.2 apresenta as estatísticas básicas em relação ao componente da prova que avalia a Formação Geral dos estudantes concluintes.
Os alunos de todo Brasil obtiveram desempenho médio de 49,7.
Quanto à variabilidade, o desvio padrão das notas dos estudantes do Brasil como um todo foi 18,2.
A maior média foi obtida na região Norte ( 52,0 ), e a menor, na região Sudeste ( 48,9 ).
As demais médias foram : 51,2 na região Nordeste, 50,5 na região Sul e 49,2 na região Centro-Oeste.
Já o maior desvio padrão foi obtido na região Nordeste ( 18,6 ) e o menor na região Sul ( 17,3 ).
Os demais desvios padrões foram : 17,9 na região Norte, 18,2 na região Sudeste e 18,5 na região Centro-Oeste.
A maior nota no componente de Formação Geral da prova do ENADE foi obtida por pelo menos um aluno da região Nordeste ( 99,0 ) enquanto que a menor nota máxima foi obtida na região Norte ( 88,5 ).
Nas outras regiões as notas máximas foram : 98,0 na região Sudeste, 92,5 na região Sul e 94,0 na região Centro-Oeste.
A mediana do Brasil como um todo foi 51,5, sendo a menor mediana encontrada nas regiões Sudeste e Centro-Oeste O Gráfico 3.5 propicia a avaliação do desempenho dos estudantes no componente de Formação Geral a partir do histograma da distribuição das notas correspondentes ao Bacharelado ( azul ), Licenciatura ( verde ), Sistemas de Informação ( bege ) e Engenharia de Computação ( roxo ).
As distribuições são unimodais, com modas no mesmo intervalo, ( 50 ; 60 ], enquanto na prova como um todo ( Gráfico 3.1 ), a moda foi alcançada em intervalos inferiores : ( 20 ; 30 ] para Bacharelado e Sistemas de Informação e ( 30 ; 40 ] para Licenciatura e Engenharia de Computação.
Nota -se, ainda, que no Gráfico 3.5 as notas apresentam uma maior dispersão do que no gráfico da distribuição das notas da prova, confirmado pela comparação dos desvios padrões : 11,6 para a nota da prova como um todo e 18,2 para o componente de Formação Geral.
Para o componente de Formação Geral, o coeficiente de assimetria da distribuição das notas dos estudantes, diferentemente do que ocorreu com a prova como um todo, é negativo para as quatro habilitações.
Observa -se que as distribuições têm causa mais longa à esquerda, nos intervalos abaixo das modas.
No entanto, nota -se um deslocamento dos histogramas : para as distribuições de Engenharia de Computação e Licenciatura a proporção de notas nos intervalos mais altos do que para as outras duas habilitações.
Os coeficientes de assimetria são : – 0,51 para o Bacharelado, – 0,63 para Licenciatura, – 0,35 para Sistemas de Informação e – 0,57 para Engenharia de Computação.
Em todas as Grandes Regiões os histogramas também possuem assimetria negativa, tomando -se os alunos em conjunto ou por habilitação.
A assimetria à esquerda mais acentuada é a da Licenciatura da região Norte ( – 1,91 ) e a mais suave é de alunos também de Licenciatura, mas da região Centro-Oeste ( – 0,16 ). ( 50,5 ) e a maior encontrada nas regiões Norte e Nordeste ( 53,5 ).
A nota mínima nesta parte foi zero em todas as regiões, sem exceção.
O Gráfico 3.5 propicia a avaliação do desempenho dos estudantes no componente de Formação Geral a partir do histograma da distribuição das notas correspondentes ao Bacharelado ( azul ), Licenciatura ( verde ), Sistemas de Informação ( bege ) e Engenharia de Computação ( roxo ).
As distribuições são unimodais, com modas no mesmo intervalo, ( 50 ; 60 ], enquanto na prova como um todo ( Gráfico 3.1 ), a moda foi alcançada em intervalos inferiores : ( 20 ; 30 ] para Bacharelado e Sistemas de Informação e ( 30 ; 40 ] para Licenciatura e Engenharia de Computação.
Nota -se, ainda, que no Gráfico 3.5 as notas apresentam uma maior dispersão do que no gráfico da distribuição das notas da prova, confirmado pela comparação dos desvios padrões : 11,6 para a nota da prova como um todo e 18,2 para o componente de Formação Geral.
Para o componente de Formação Geral, o coeficiente de assimetria da distribuição das notas dos estudantes, diferentemente do que ocorreu com a prova como um todo, é negativo para as quatro habilitações.
Observa -se que as distribuições têm causa mais longa à esquerda, nos intervalos abaixo das modas.
No entanto, nota -se um deslocamento dos histogramas : para as distribuições de Engenharia de Computação e Licenciatura a proporção de notas nos intervalos mais altos do que para as outras duas habilitações.
Os coeficientes de assimetria são : – 0,51 para o Bacharelado, – 0,63 para Licenciatura, – 0,35 para Sistemas de Informação e – 0,57 para Engenharia de Computação.
Em todas as Grandes Regiões os histogramas também possuem assimetria negativa, tomando -se os alunos em conjunto ou por habilitação.
A assimetria à esquerda mais acentuada é a da Licenciatura da região Norte ( – 1,91 ) e a mais suave é de alunos também de Licenciatura, mas da região Centro-Oeste ( – 0,16 ).
Nos Gráficos 3.6, 3.7 e 3.8 são apresentadas as informações referentes ao desempenho dos Concluintes no componente de Formação Geral, em diferentes agregações : Grande Região do país, Categoria Administrativa e Organização Acadêmica.
Observa -se pelo Gráfico 3.6 que, no Componente de Formação Geral, existe diferença estatisticamente significativa entre as duas maiores médias, obtidas nas regiões Norte ( 52,0 ) e Nordeste ( 51,2 ) das notas, e as duas menores, obtidas nas regiões Sudeste ( 48,9 ) e Centro-Oeste ( 49,2 ).
No Gráfico 3.7, que representa as notas médias no Componente de Formação Geral segundo Categoria Administrativa do país, observa -se que existe diferença estatisticamente significativa entre as médias.
Como para as notas da prova como um todo, no componente de Formação Geral, os concluintes de Computação ( Bacharelado, Licenciatura, Sistemas de Informação e Engenharia de Computação, considerados em conjunto ), das IES Públicas ( 51,1 ) obtiveram uma média maior do que os das IES Privadas ( 49,3 ).
Considerando -se o tipo de Organização Acadêmica, nota -se, no Gráfico 3.8, que não há diferença estatisticamente significativa entre as médias dos concluintes de Universidades ( 50,1 ), Centros Universitários e Faculdades ( ambas iguais a 49,3 ). 3.1.3 Estatísticas Básicas do Componente de Conhecimento Específico A Tabela 3.3 apresenta as estatísticas básicas referentes ao componente de Conhecimento Específico da área de Computação ( Bacharelado, Licenciatura, Sistemas de Informação e Engenharia de Computação ).
A média do desempenho dos alunos do Brasil como um todo foi 25,8.
A maior média foi obtida na região Sul ( 27,7 ), e a menor, na região Norte ( 24,0 ).
As demais médias foram : 27,4 na região Nordeste ; 25,0 na região Sudeste ; e 24,6 na região Centro-Oeste.
Quanto à variabilidade das notas, o desvio padrão do Brasil como um todo foi 12,7, sendo o maior desvio padrão observado na região Nordeste ( 13,5 ) e o menor na região Norte ( 11,6 ).
Os demais desvios foram : 12,4 da região Sudeste ; 13,2 da região Sul ; e 12,3 da região Centro-Oeste.
A mediana das notas dos estudantes de todo o Brasil foi 24,1.
A maior mediana ocorreu na região Sul ( 25,3 ) e a menor nas regiões Norte, Sudeste e Centro-Oeste ( 23,6 ).
A nota máxima do Brasil como um todo foi 93,7, sendo obtida por pelo menos um aluno da região Sudeste.
As demais notas máximas foram : 74,0 na região Norte ; 83,3 na região Nordeste ; 86,2 na região Sul ; e 78,6 na região Centro-Oeste.
A nota mínima foi zero em todas as regiões.
Assim como os Gráficos 3.1 e 3.5, o Gráfico 3.9, apresentado a seguir, proporciona uma avaliação do desempenho de concluintes dos cursos da Área de Computação ( Bacharelado, Licenciatura, Sistemas da Informação e Engenharia de Computação ) em relação ao componente de Conhecimento Específico com histogramas das distribuições das notas correspondentes.
Dentre as distribuições apresentadas até aqui, esta é a mais concentrada em notas baixas.
Estas também são distribuições unimodais e o grupo modal é o ( 20 ; 30 ] para todos as habilitações : Bacharelado, Licenciatura, Sistemas de Informação e Engenharia de Computação.
Como nos gráficos anteriores a distribuição de Engenharia de Computação tem frequências maiores dos que as das outras habilitações nos intervalos de notas mais altas.
O coeficiente de assimetria das quatro distribuições do componente de Conhecimento Específico é positivo : 0,77 para a habilitação de Bacharelado, 0,45 para a Licenciatura, 0,50 para Sistemas de Informação e 0,71 para Engenharia de Computação.
Nota -se pelo histograma ( Gráfico 3.9 ) que há cauda mais longa à esquerda para todos as habilitações.
Apenas para a Licenciatura na região Norte o coeficiente de assimetria é negativo.
Para todos as outras habilitações e regiões há assimetria à direita, coeficiente de assimetria positivo, sendo que a mais acentuada é a distribuição de notas dos alunos de Bacharelado da região Sudeste ( 0,85 ) e a mais leve encontrada na região Centro-Oeste para alunos de Sistemas de Informação ( 0,32 ).
Os Gráficos 3.10, 3.11 e 3.12 apresentam uma comparação dos resultados em relação à Grande Região do país, à Categoria Administrativa e à Organização Acadêmica, agora levando em conta o desempenho de alunos no componente de Conhecimento Específico.
Pelo Gráfico 3.10, observa -se que não existe diferença estatisticamente significativa entre as duas maiores médias que ocorreram nas regiões Sul ( 27,7 ) e Nordeste ( 27,4 ).
Já entre estas e as demais médias, das regiões Centro-Oeste ( 24,6 ), Sudeste ( 25,0 ) e Norte ( 24,0 ) há diferença estatisticamente significativa ao nível de 95 \%.
Quanto à Categoria Administrativa ( Gráfico 3.11 ), observa -se um comportamento semelhante àquele da parte de Formação Geral e à prova como um todo, ou seja, existe diferença estatisticamente significativa entre as médias das IES Públicas ( 31,9 ) e Privadas ( 23,9 ), sendo que neste caso, em relação ao componente de Conhecimento Específico, esta diferença é maior do que nos casos anteriores, prova como um todo e Componente de Formação Geral.
Quanto ao Gráfico 3.12, observa -se, mais uma vez, que existe diferença estatisticamente significativa ao nível de 95 \% entre as notas no componente de Conhecimento Específico dos diferentes tipos de Organização Acadêmica.
Sendo que a média dos Concluintes das Universidades ( 28,1 ) foi maior do que de Centros Universitários ( 23,5 ), e esta maior do que a de Faculdades ( 22,7 ). 3.2 ANÁLISE DAS QUESTÕES OBJETIVAS 3.2.1 Componente de Formação Geral A Tabela 3.4 apresenta as estatísticas básicas relativas às oito questões objetivas do componente da prova que abrange a Formação Geral dos estudantes.
A média do Brasil foi 50,6.
A menor média foi encontrada na região Sudeste ( 49,4 ) e a maior na região Nordeste ( 52,7 ).
As demais médias foram : 52,5 na região Norte, 51,4 na região Sul e 50,7 na região Centro-Oeste.
O desvio padrão do Brasil foi 19,4, sendo o maior desvio padrão encontrado na região Centro-Oeste ( 20,0 ) e o menor na região Sul ( 18,6 ).
Os demais desvios foram : 19,5 nas regiões Norte e Nordeste e 19,4 na região Sudeste.
As medianas ( 50,0 ), as notas máximas ( 100,0 ) e as notas mínimas ( 0,0 ) foram iguais para todas as regiões.
A Tabela 3.5 apresenta os índices de facilidade e discriminação ( ponto bisserial ) das questões objetivas do componente de Formação Geral.
Quanto ao índice de facilidade, foram usadas as seguintes cores para diferenciar o nível de dificuldade da questão : Azul para as questões classificadas com índice muito fácil ( índice \&gt ; =0,86 ), verde para as questões classificadas com índice fácil ( 0,61 a 0,85 ), amarelo para as questões classificadas com médio ( 0,41 a 0,60 ), vermelho para as questões classificadas com difícil ( 0,16 a 0,40 ) e roxo para as questões classificadas com muito difícil ( \&lt ; =0,15 ).
Já quanto ao índice de discriminação, foram usadas as seguintes cores para qualificar a questão : As questões classificadas com índice fraco receberam a cor vermelho ( índice \&lt ; =0,19 ), as classificadas com médio receberam a cor amarelo ( 0,20 a 0,29 ), as classificadas com bom receberam a cor verde ( 0,30 a 0,39 ) e as classificadas com muito bom ( \&gt ; =0,40 ) receberam a cor azul.
As questões objetivas do componente de Formação Geral, segundo o índice de facilidade, foram assim avaliadas : das oito questões, nenhuma teve o índice de facilidade classificado como muito fácil.
Quatro questões foram tidas como fáceis, por terem índice de acertos situado na faixa entre 0,61 e 0,85 ( de 61,0 \% a 85,0 \% de acertos ).
Uma questão foi considerada de dificuldade média, situando -se no intervalo entre 0,41 e 0,60 do índice de facilidade ( entre 41,0 \% e 60,0 \% acertos ).
Outras duas questões foram classificadas na categoria difícil, situando -se no intervalo entre 0,16 e 0,40.
Por fim, uma das questões apresentou menos de 15 \% de acertos, razão pela qual foi classificada como muito difícil.
Como já comentado, para análise das questões objetivas relativas à Formação Geral segundo o poder de discriminação, utilizou -se, o índice de discriminação ponto bisserial.
Nesta análise as questões foram assim avaliadas : quatro das oito questões apresentaram A Tabela 3.5 apresenta os índices de facilidade e discriminação ( ponto bisserial ) das questões objetivas do componente de Formação Geral.
Quanto ao índice de facilidade, foram usadas as seguintes cores para diferenciar o nível de dificuldade da questão : Azul para as questões classificadas com índice muito fácil ( índice \&gt ; =0,86 ), verde para as questões classificadas com índice fácil ( 0,61 a 0,85 ), amarelo para as questões classificadas com médio ( 0,41 a 0,60 ), vermelho para as questões classificadas com difícil ( 0,16 a 0,40 ) e roxo para as questões classificadas com muito difícil ( \&lt ; =0,15 ).
Já quanto ao índice de discriminação, foram usadas as seguintes cores para qualificar a questão : As questões classificadas com índice fraco receberam a cor vermelho ( índice \&lt ; =0,19 ), as classificadas com médio receberam a cor amarelo ( 0,20 a 0,29 ), as classificadas com bom receberam a cor verde ( 0,30 a 0,39 ) e as classificadas com muito bom ( \&gt ; =0,40 ) receberam a cor azul.
As questões objetivas do componente de Formação Geral, segundo o índice de facilidade, foram assim avaliadas : das oito questões, nenhuma teve o índice de facilidade classificado como muito fácil.
Quatro questões foram tidas como fáceis, por terem índice de acertos situado na faixa entre 0,61 e 0,85 ( de 61,0 \% a 85,0 \% de acertos ).
Uma questão foi considerada de dificuldade média, situando -se no intervalo entre 0,41 e 0,60 do índice de facilidade ( entre 41,0 \% e 60,0 \% acertos ).
Outras duas questões foram classificadas na categoria difícil, situando -se no intervalo entre 0,16 e 0,40.
Por fim, uma das questões apresentou menos de 15 \% de acertos, razão pela qual foi classificada como muito difícil.
Como já comentado, para análise das questões objetivas relativas à Formação Geral segundo o poder de discriminação, utilizou -se, o índice de discriminação ponto bisserial.
Nesta análise as questões foram assim avaliadas : quatro das oito questões apresentaram índices acima ou igual a 0,40 e, assim, foram classificadas com índice muito bom para esse grupo de estudantes.
Três questões tiveram bom índice de discriminação, com valor entre 0,30 e 0,39, e uma foi classificada com índice médio para esse grupo de estudantes.
Nenhuma questão teve nível fraco de discriminação.
O índice de facilidade variou de 0,15 a 0,72, e o de discriminação, de 0,26 a 0,56.
As quatro questões com índices de discriminação muito bom, as de números 1, 2, 3 e 5, figuraram entre dois níveis de dificuldade : três classificadas na categoria fácil ( questões 1, 3 e 5 ) do índice de facilidade, uma na categoria difícil ( questão 2 ).
Em particular, a questão 3 foi a que apresentou o maior poder discriminatório, com índice 0,56, e foi também uma das mais fáceis, com uma proporção de 0,72 acertos, o máximo de acertos.
Esta proporção de acertos, 72,0 \%, também foi alcançada pela questão 6.
A questão de número 7 apresentou um índice de facilidade de 0,28, ou seja, apenas 28 \% dos estudantes conseguiram resolvê -la, dentro do universo de participantes, e seu índice de discriminação foi bom ( 0,35 ).
Já a questão 8, que obteve índice de discriminação médio, 0,26, obteve índice de facilidade 0,15 e foi classificada como muito difícil.
O Gráfico 3.13, para exemplificar, analisa o comportamento da questão de número 3 de Formação Geral.
Trata -se de uma das duas questões mais fáceis e a que obteve o maior índice de discriminação dessa parte da prova.
Neste gráfico, cada uma das cinco curvas representa o percentual de respostas em determinada alternativa da questão, em função do número de acertos dos estudantes nesta parte da prova ( Formação Geral/Múltipla Escolha ), antes de possíveis eliminações pelo critério do ponto bisserial.
Em princípio, a soma das escolhas possíveis deveria ser igual a 100 \%.
Não é este o caso, pois são computadas no gráfico a quantidade de questões deixadas em branco ou com múltiplas respostas.
Como foram oito as questões, os valores variam de 0 a 8 acertos.
A curva em vermelho corresponde à alternativa E, a correta para esta questão.
Assim, observa -se que entre os estudantes com menor número de acertos, nessa parte do exame, a situação mais frequente foi a escolha de uma das alternativas incorretas : a alternativa A ( em azul ).
Na medida em que o número de acertos aumenta, indicando desempenho melhor nesta parte da prova, aumenta concomitantemente a proporção de estudantes que selecionaram a alternativa correta E, atingindo 100 \% para os estudantes com 8 acertos.
Essa análise permite verificar como a questão discriminou os grupos de desempenho, justificando o alto índice obtido na questão.
Os gráficos relativos às demais questões de Formação Geral constam do Anexo I. 3.2.2 Componente de Conhecimento Específico – Computação ( Bacharelado – Ciência da Computação ) A Tabela 3.6 apresenta as estatísticas básicas em relação às questões objetivas do componente de Conhecimento Específico da prova de Computação – Bacharelado, por Grande Região.
A média do Brasil deste componente foi de 30,8.
A menor média foi observada na região Norte ( 27,4 ) e a maior na região Sul ( 34,2 ).
O desvio padrão de todo o Brasil foi 15,4, sendo o menor desvio padrão encontrado na região Norte ( 13,9 ) e o maior na região Nordeste ( 16,6 ).
A mediana de todo o Brasil foi 28,6, a mesma encontrada na região Sudeste e na Centro-Oeste.
A maior mediana foi 33,3, das regiões Nordeste e Sul, e a menor mediana ocorreu na região Norte, 23,8.
A nota máxima foi 90,5, obtida nas questões objetivas do componente de Conhecimento Específico, por pelo menos um aluno das regiões Sudeste e Sul.
A menor nota máxima foi a da região Norte, 76,2.
A Tabela 3.7 apresenta os índices de facilidade e discriminação ( ponto bisserial ) das questões objetivas do componente de Conhecimento Específico para os estudantes de Computação – Bacharelado.
Para facilitar a diferenciação das questões usaram -se as mesmas cores das Tabelas 3.5 para as diferentes classificações dos índices de facilidade e de discriminação.
Dentre as 27 questões objetivas da parte da prova relativa ao componente de Conhecimento Específico para os concluintes do Bacharelado ( questões de 9 a 30 e de 36 a 40 ), a questão 13 foi anulada pela Comissão.
Desse modo, a classificação quanto ao índice de facilidade foi estabelecida com base em 26 das 27 questões.
A partir dos índices obtidos, pode -se concluir que a maior parte das questões objetivas da prova foi considerada pelo menos difícil : das 26 questões válidas, 18 foram classificadas como difíceis ( a classificação modal ) e três como muito difíceis.
Não houve questão classificada como fácil ou muito fácil, ao passo que cinco foram tidas como médias, na faixa de 0,41 a 0,60 do índice de facilidade.
Já quanto aos índices de discriminação das questões objetivas do componente de Conhecimento Específico da prova, tem -se como resultado a seguinte classificação : nove das 26 questões foram consideradas como boas, enquanto quatro delas tiveram índice de discriminação muito bom.
Assim, para a metade das questões – 13 em 26 – os índices de discriminação foram bons ou muito bons.
Dentre as demais, oito delas foram classificadas como médias e outras cinco como fracas, sendo também 13 a quantidade de questões nos dois patamares mais baixos de discriminação.
Constata -se, assim, que a prova – no que se refere ao componente de Conhecimento Específico – possuía capacidade média de discriminar entre aqueles que dominam ou não o conteúdo.
Dentre as questões que alcançaram os maiores índices de discriminação, quatro delas, as de números 12, 14, 19 e 21 foram classificadas com índice de discriminação muito bom, situando -se no intervalo de 0,40 a 0,50 do índice, três delas ( questões 12, 19 e 21 ) foram classificadas na categoria médio, quanto ao índice de facilidade, e a questão 14 na categoria difícil.
A questão de número 38 foi a mais difícil dentre as 26 questões específicas válidas, com baixo índice de facilidade, apenas 8,0 \% de acertos.
Essa questão apresentou poder discriminatório igualmente baixo, 0,15, o que comprova ter sido esta a mais difícil para os estudantes.
Destaca -se, também, a questão 40, com índice de facilidade 0,09, também considerada muito difícil, já que apenas 9,0 \% dos estudantes responderam esta questão acertadamente.
O índice de discriminação foi também baixo, 0,14.
Tais questões foram, portanto, pelo critério ponto bisserial, consideradas inadequadas.
Por isso, as questões 38 e 40 foram eliminadas do cômputo da nota final.
Além destas duas, as demais questões com índice fraco de discriminação, questões 17, 23 e 24 também não foram utilizadas no cômputo final das notas, num total de cinco questões eliminadas.
A título de exemplo das análises do comportamento das questões objetivas, o Gráfico 3.14 analisa a questão 12 do componente de Conhecimento Específico.
Esta foi a questão com índice de discriminação mais alto da prova ( 0,51 ), muito bom.
Esta questão apresentou índice de facilidade 0,41, ou seja, 41,0 \% dos estudantes assinalaram acertadamente a opção B, correspondente ao gabarito.
Assim, seu índice de facilidade foi classificado como médio.
Neste gráfico, cada uma das cinco curvas representa o percentual de respostas em determinada alternativa da questão 12, em função do número de acerto dos estudantes nesta parte da prova, antes de possíveis eliminações de questões pelo critério do ponto bisserial.
A alternativa correta B, representada no gráfico pela curva em verde, foi escolhida em maiores proporções pelos alunos com desempenho melhor nesta parte da prova.
Já as alternativas incorretas, também denominadas distratores, foram selecionadas principalmente por aqueles com notas mais baixas.
Neste caso também a soma não é sempre 100 \% por causa das questões não respondidas ou com mais de uma opção marcada.
Aqueles com nota zero, na sua quase totalidade deixaram esta questão em branco ou marcaram mais de uma alternativa, comportamento considerado inválido.
A proporção de alunos que selecionou a resposta correta B aumenta gradativamente, chegando a atingir 100 \% para mais de 20 acertos, enquanto a proporção dos que escolheram alternativas incorretas decai, a partir de cinco acertos, como função do número de acertos nesta parte da prova.
Os gráficos relativos às demais questões do componente de Conhecimento Específico constam do Anexo I. 3.2.3 Componente de Conhecimento Específico – Computação ( Licenciatura ) A Tabela 3.8 apresenta as estatísticas básicas em relação às questões objetivas do componente de Conhecimento Específico da prova de Computação - Licenciatura, por Grande Região.
A média do Brasil deste componente foi de 29,2.
A maior média foi observada na região Norte ( 32,3 ) e a menor na região Sul ( 23,9 ).
O desvio padrão de todo o Brasil foi 13,6, sendo o maior desvio padrão encontrado na região Sul ( 14,6 ) e o menor na região Norte ( 11,3 ).
A mediana de todo o Brasil foi 27,8, também encontrada para as regiões Sudeste, Sul e Centro-Oeste.
A maior mediana ocorreu nas regiões Norte e Nordeste ( 33,3 ).
A nota máxima da prova foi 72,2, obtida por pelo menos um concluinte de Licenciatura das regiões Nordeste, Sudeste e Centro-Oeste, nas questões objetivas do componente de Conhecimento Específico.
A menor nota máxima ocorreu na região Norte ( 55,6 ).
A Tabela 3.9 apresenta os índices de facilidade e discriminação ( ponto bisserial ) das questões objetivas do componente de Conhecimento Específico para os estudantes de Computação – Licenciatura.
Para facilitar a diferenciação das questões usaram -se as mesmas cores das Tabelas 3.5 para as diferentes classificações dos índices de facilidade e de discriminação.
Dentre as questões objetivas da parte da prova relativa ao componente de Conhecimento Específico para os alunos de Licenciatura ( questões de 9 a 35 ), a questão 13 foi anulada pela Comissão.
Desse modo, a classificação quanto ao índice de facilidade foi estabelecida com base em 26 das 27 questões.
A partir dos índices obtidos, pode -se concluir que a maior parte das questões objetivas da prova foi considerada pelo menos difícil : das 26 questões válidas, 17 foram classificadas como difíceis ( a classificação modal ) e cinco como muito difíceis.
Não houve questão classificada como muito fácil, ao passo que uma foi tida como fácil, na faixa de 0,61 a 0,85 do índice de facilidade, e outras três consideradas médias, entre 0,41 e 0,60.
Já quanto aos índices de discriminação das questões objetivas do componente de Conhecimento Específico da prova, tem -se como resultado a seguinte classificação : nove das 26 questões foram consideradas como boas, enquanto apenas uma delas teve índice de discriminação muito bom.
Assim, para menos da metade das questões – 10 em 26 – os índices de discriminação foram bons ou muito bons.
Dentre as demais, oito delas foram classificadas como médias e outras oito como fracas, sendo dezesseis, por conseguinte, a quantidade de questões nos dois patamares mais baixos de discriminação.
Constata -se, assim, que a prova – no que se refere ao componente de Conhecimento Específico – possuía capacidade abaixo da média de discriminar entre aqueles que dominam ou não o conteúdo.
A questão que alcançou o maior índice de discriminação, a de número 19, teve índice 0,40, no limite inferior do intervalo que classifica a discriminação como muito boa, de 0,40 a 0,50.
Esta questão foi classificada na categoria difícil, quanto ao índice de facilidade.
As questões de números 11, 15, 16, 26 e 29 foram classificadas como muito difíceis, sendo que dentre elas a de número 29 foi a que obteve menor número de acertos, 8,0 \% dos alunos responderam esta questão acertadamente.
Dentre as mais difíceis, três delas tiveram índice de discriminação fraco, as questões 11, 15 e 26, o que comprova terem sido, as mais difíceis para os estudantes.
Destacam -se, também, as questões 9, 17, 20, 23 e 31, classificadas como difíceis, com índices de facilidade entre 0,21 e 0,40.
Tais questões tiveram índice de discriminação fraco, abaixo de 0,20.
Tais questões foram, portanto, pelo critério ponto bisserial consideradas inadequadas.
Por isso, as questões 9, 11, 15, 17, 20, 23, 26 e 31 foram eliminadas do cômputo da nota final, num total de oito questões eliminadas.
A título de exemplo das análises do comportamento das questões objetivas, o Gráfico 3.15 analisa a questão 19 do componente de Conhecimento Específico.
Esta foi considerada uma questão de dificuldade média, apresentando índice de facilidade 0,33, ou seja, 33,0 \% dos estudantes assinalaram acertadamente a opção E, correspondente ao gabarito.
Seu índice de discriminação foi o mais alto alcançado na prova para este grupo de estudantes, igual a 0,40, classificado como muito bom.
Neste gráfico, cada uma das cinco curvas representa o percentual de respostas em determinada alternativa da questão 19, em função do número de acerto dos estudantes nesta parte da prova, antes de possíveis eliminações de questões pelo critério do ponto bisserial.
A alternativa correta E, representada no gráfico pela curva em vermelho, foi escolhida em maiores proporções pelos alunos com desempenho melhor nesta parte da prova e, por isso, não aparece esta abscissa no eixo horizontal.
Já as alternativas incorretas, também denominadas distratores, foram selecionadas principalmente por aqueles com notas mais baixas.
Neste caso também a soma não é sempre 100 \% por causa das questões não respondidas ou com mais de uma opção marcada.
Aqueles com nota zero, na sua quase totalidade deixaram esta questão em branco ou marcaram mais de uma alternativa, comportamento considerado inválido.
A proporção de alunos que selecionou a resposta correta E aumenta gradativamente, chegando a atingir 100 \% para 14 e mais de 16 acertos, enquanto a proporção dos que escolheram alternativas incorretas decai, a partir de três ou quatro acertos, como função do número de acertos nesta parte da prova.
Os gráficos relativos às demais questões do componente de Conhecimento Específico constam do Anexo I. 3.2.4 Componente de Conhecimento Específico – Computação ( Sistemas de Informação ) A Tabela 3.10 apresenta as estatísticas básicas em relação às questões objetivas do componente de Conhecimento Específico da prova de Computação – Sistemas de Informação, por Grande Região.
A média do Brasil deste componente foi de 35,1.
A menor média foi observada na região Centro-Oeste ( 32,9 ) e a maior na região Nordeste ( 39,9 ).
O desvio padrão de todo o Brasil foi 14,7, sendo o maior desvio padrão encontrado na região Nordeste ( 15,8 ) e o menor na região Norte ( 12,8 ).
A mediana de todo o Brasil foi 33,3, também encontrada para as regiões Sudeste e Centro-Oeste.
As medianas das regiões Norte, Nordeste e Sul foram maiores e iguais, 38,1.
A nota máxima da prova foi 95,2, obtida por pelo menos um concluinte de Sistemas da Informação da região Sudeste, nas questões objetivas do componente de Conhecimento Específico.
A menor nota máxima foi obtida na região Norte, 66,7.
A Tabela 3.11 apresenta os índices de facilidade e discriminação ( ponto bisserial ) das questões objetivas do componente de Conhecimento Específico para os estudantes de Computação – Sistemas de Informação.
Para facilitar a diferenciação das questões usaram -se as mesmas cores das Tabelas 3.5 para as diferentes classificações dos índices de facilidade e de discriminação.
Dentre as questões objetivas da parte da prova relativa ao componente de Conhecimento Específico para os alunos de Sistemas de Informação ( questões de 9 a 30 e de 46 a 50 ), a questão 13 foi anulada pela Comissão.
Desse modo, a classificação quanto ao índice de facilidade foi estabelecida com base em 26 das 27 questões.
A partir dos índices obtidos, pode -se concluir que a maior parte das questões objetivas da prova foi considerada pelo menos difícil : das 26 questões válidas, 17 foram classificadas como difíceis ( a classificação modal ) e cinco como muito difíceis.
Não houve questão classificada como muito fácil, ao passo que uma foi tida como fácil, na faixa de 0,61 a 0,85 do índice de facilidade, e outras três consideradas médias, entre 0,41 e 0,60.
Já quanto aos índices de discriminação das questões objetivas do componente de Conhecimento Específico da prova, tem -se como resultado a seguinte classificação : quatro das 26 questões foram consideradas como boas, enquanto nenhuma questão obteve índice de discriminação muito bom.
Assim, para a maioria das questões – 22 em 26 – os índices de discriminação foram médios ou fracos, os dois patamares mais baixos de discriminação.
Dentre elas, 14 foram classificadas como médias e outras oito como fracas.
Constata -se, assim, que a prova – no que se refere ao componente de Conhecimento Específico – possuía capacidade abaixo da média de discriminar entre aqueles que dominam ou não o conteúdo.
Dentre as questões que alcançaram índices de discriminação bom, situando -se no intervalo de 0,40 a 0,50 do índice, quatro delas, as de números 18, 19, 21 e 48, uma delas ( questão 48 ) foi classificada na categoria fácil, a única questão da prova neste grau de dificuldade.
As questões 18 e 19 ficaram na categoria média de dificuldade e a questão 21 foi considerada difícil.
A questão de número 29 foi a mais difícil dentre as 26 questões específicas válidas, com baixo índice de facilidade, apenas 7,0 \% de acertos.
Essa questão apresentou poder discriminatório médio, 0,21, o que mostra que apesar de muito difícil ela teve capacidade média de discriminar os estudantes.
Destacam -se, também, as questões 11 e 26 que tiveram índices de facilidade muito difícil, 0,10 e 0,15, respectivamente, ou seja, na questão 11 apenas 10,0 \% dos estudantes responderam acertadamente e na questão 26 foram 15,0 \% os que acertaram.
Tais questões tiveram índices de discriminação igualmente baixos, 0,15 para a questão 11 e 0,14 para a questão 26.
Portanto, pelo critério ponto bisserial consideradas inadequadas.
Por isso, as questões 11 e 26 foram eliminadas do cômputo da nota final.
Além destas duas, as demais questões com índice fraco de discriminação, questões 9, 12, 17, 23, 24 e 47, também não foram utilizadas no cômputo final das notas, num total de oito questões eliminadas.
A título de exemplo das análises do comportamento das questões objetivas, o Gráfico 3.16 analisa a questão 19 do componente de Conhecimento Específico.
Esta foi considerada uma questão de dificuldade média, apresentando índice de facilidade 0,44, ou seja, 44,0 \% dos estudantes assinalaram acertadamente a opção E, correspondente ao gabarito.
Seu índice de discriminação foi igual a 0,39, classificado como bom e, apesar disso, foi a questão que apresentou o maior índice de discriminação.
Neste gráfico, cada uma das cinco curvas representa o percentual de respostas em determinada alternativa da questão 19, em função do número de acerto dos estudantes nesta parte da prova, antes de possíveis eliminações de questões pelo critério do ponto bisserial.
A alternativa correta E, representada no gráfico pela curva em vermelha, foi escolhida em maiores proporções pelos alunos com desempenho melhor nesta parte da prova.
Já as alternativas incorretas, também denominadas distratores, foram selecionadas principalmente por aqueles com notas mais baixas.
Neste caso também a soma não é sempre 100 \% por causa das questões não respondidas ou com mais de uma opção marcada.
Aqueles com nota zero, na sua quase totalidade deixaram esta questão em branco ou marcaram mais de uma alternativa, comportamento considerado inválido e, por isso, não aparece esta abscissa no eixo horizontal.
A proporção de alunos que selecionou a resposta correta E aumenta gradativamente, chegando a atingir 100 \% para os que obtiveram 16, 19 e 20 acertos, enquanto a proporção dos que escolheram alternativas incorretas decai, a partir de dois acertos, como função do número de acertos nesta parte da prova.
Os gráficos relativos às demais questões do componente de Conhecimento Específico constam do Anexo I. 3.2.5 Componente de Conhecimento Específico – Computação ( Engenharia de Computação ) A Tabela 3.12 apresenta as estatísticas básicas em relação às questões objetivas do componente de Conhecimento Específico da prova de Computação – Engenharia de Computação, por Grande Região.
A média do Brasil deste componente foi de 27,9.
A menor média foi observada na região Centro-Oeste ( 26,5 ) e a maior na região Sul ( 29,7 ).
O desvio padrão de todo o Brasil foi 12,4, sendo o maior desvio padrão encontrado na região Sul ( 13,0 ) e o menor na região Norte ( 11,4 ).
A mediana de todo o Brasil foi 27,8, sendo a mesma para todas as regiões.
A nota máxima da prova foi 94,4, obtida por pelo menos um concluinte de Engenharia de Computação da região Sul, nas questões objetivas do componente de Conhecimento Específico.
A menor nota máxima ( 66,7 ) ocorreu nas regiões Norte e Centro-Oeste.
A Tabela 3.13 apresenta os índices de facilidade e discriminação ( ponto bisserial ) das questões objetivas do componente de Conhecimento Específico.
Para facilitar a diferenciação das questões usaram -se as mesmas cores da Tabela 3.5 para as diferentes classificações dos índices de facilidade e de discriminação.
Dentre as questões objetivas da parte da prova relativa ao componente de Conhecimento Específico para alunos de Engenharia da Computação ( questões 9 a 30 e 41 a 45 ), duas questões foram anuladas pela Comissão : questões 13 e 44.
Desse modo, a classificação quanto ao índice de facilidade foi estabelecida com base em 25 das 27 questões.
A partir dos índices obtidos, pode -se concluir que a maioria das questões objetivas da prova foi considerada difícil : 18 das 25 questões.
Não houve questão classificada como muito difícil nem muito fácil, ao passo que duas foram tidas como fáceis, na faixa de 0,61 a 0,85 do índice de facilidade, e cinco consideradas médias, entre 0,41 e 0,60.
Já quanto aos índices de discriminação das questões objetivas do componente de Conhecimento Específico da prova, tem -se como resultado a seguinte classificação : quatro das 26 questões foram consideradas como muito boas, enquanto nove delas tiveram índice de discriminação bom.
Assim, para a maior parte das questões – 13 em 25 – os índices de discriminação foram bons ou muito bons.
Dentre as demais, oito delas foram classificadas como médias e outras quatro como fracas, sendo doze, por conseguinte, a quantidade de questões nos dois patamares mais baixos de discriminação.
Constata -se, assim, que a prova – no que se refere ao componente de Conhecimento Específico – possuía boa capacidade de discriminar entre aqueles que dominam ou não o conteúdo.
Todas as questões que alcançaram os maiores índices de discriminação, as de números 12, 14, 19 e 21, classificadas com índice muito bom, situando -se no intervalo de 0,42 a 0,50 do índice, quanto ao índice de facilidade, foram classificadas como sendo de dificuldade média.
Já a questão de número 11, a mais difícil dentre as 25 questões específicas válidas, teve índice de discriminação bom ( 0,30 ).
Dentre as 18 questões consideradas difíceis estão incluídas as quatro questões com índice de discriminação fraco, as de números 17, 23, 24 e 41, que foram eliminadas do cômputo da nota final pelo critério do índice de discriminação ponto bisserial.
A título de exemplo das análises do comportamento das questões objetivas, o Gráfico 3.17 analisa a questão 19 do componente de Conhecimento Específico.
Esta questão foi respondida acertadamente pela mais da metade dos estudantes apresentando índice de facilidade médio ( 0,58 ), já que 58,0 \% dos estudantes assinalaram acertadamente a opção E, correspondente ao gabarito.
Seu índice de discriminação foi igual a 0,44, classificado como muito bom e o maior valor alcançado por este índice para o conjunto de perguntas em análise.
Neste gráfico, cada uma das cinco curvas representa o percentual de respostas em determinada alternativa da questão 19, em função do número de acertos dos estudantes nesta parte da prova, antes de possíveis eliminações de questões pelo critério do ponto bisserial.
A alternativa correta E, representada no gráfico pela curva em vermelho, foi escolhida em maiores proporções pelos alunos com desempenho melhor nesta parte da prova.
Já as alternativas incorretas, também denominadas distratores, foram selecionadas principalmente por aqueles com notas mais baixas.
Aqueles com nota zero, na sua quase totalidade deixaram esta questão em branco ou marcaram mais de uma alternativa, comportamento considerado inválido e, por isso, não aparece esta abscissa no eixo horizontal.
A proporção de alunos que selecionou a resposta correta E aumenta gradativamente, chegando a atingir 100 \% para alunos com 16 acertos ou mais do que 18 acertos, enquanto a proporção dos que escolheram alternativas incorretas decai, a partir de três acertos, como função do número de acertos nesta parte da prova.
Os gráficos relativos às demais questões do componente de Conhecimento Específico constam do Anexo I. 3.3 ANÁLISE DAS QUESTÕES DISCURSIVAS 3.3.1 Componente de Formação Geral As análises dos resultados de desempenho dos estudantes de Computação ( Bacharelado, Licenciatura, Sistemas de Informação e Engenharia de Computação ) nas duas questões discursivas relativas à Formação Geral encontram -se na Tabela 3.14 e no Gráfico 3.18.
Na tabela 3.14 observa -se que se por um lado a nota média nesse conjunto de questões foi um pouco menor do que a obtida nas objetivas, por outro, houve aumento no desvio padrão.
Os estudantes de todo o Brasil obtiveram, em Formação Geral, média 50,6 nas questões objetivas e 48,5 nas questões discursivas.
Já o desvio padrão passou de 19,4 nas questões objetivas do componente de Formação Geral dos alunos de todo o Brasil, para O Gráfico 3.18 representa a distribuição das notas nas questões discursivas no componente de Formação Geral segundo a opção de Bacharelado, Licenciatura, Sistemas de Informação ou Engenharia de Computação.
As distribuições são bimodais.
Para os concluintes de todos as habilitações a moda principal ocorre no intervalo ( 40 ; 50 ].
Para o Bacharelado, Sistemas de Informação e Engenharia da Computação há modas secundárias no intervalo [ 0 ; 10 ], que inclui as notas zero e os alunos que deixaram essas questões discursivas em branco.
Já para os concluintes da Licenciatura a segunda moda é o intervalo ( 70 ; 80 ].
Os coeficientes de assimetria das quatro distribuições são negativos para o Brasil como um todo : – 0,33 para o Bacharelado, – 0,53 para a Licenciatura, – 0,31 para Sistemas de Informação e – 0,43 para Engenharia de Computação.
As distribuições de todos as habilitações, em todas as regiões, têm o mesmo comportamento, têm cauda mais acentuada no lado esquerdo do gráfico, ou seja, são assimétricas à esquerda. 27,8 nas questões discursivas do mesmo componente.
A maior média foi obtida na região Norte ( 51,3 ) e a menor na região Centro-Oeste ( 47,0 ).
A mediana de todo o Brasil, neste componente, foi 50,0, a mesma de todas as Grandes Regiões, exceto na região Norte, onde a mediana foi um pouco mais alta, 55,0.
A nota máxima ( 100,0 ) e a nota mínima ( 0,0 ) também foram iguais em todas as regiões do Brasil.
O Gráfico 3.18 representa a distribuição das notas nas questões discursivas no componente de Formação Geral segundo a opção de Bacharelado, Licenciatura, Sistemas de Informação ou Engenharia de Computação.
As distribuições são bimodais.
Para os concluintes de todos as habilitações a moda principal ocorre no intervalo ( 40 ; 50 ].
Para o Bacharelado, Sistemas de Informação e Engenharia da Computação há modas secundárias no intervalo [ 0 ; 10 ], que inclui as notas zero e os alunos que deixaram essas questões discursivas em branco.
Já para os concluintes da Licenciatura a segunda moda é o intervalo ( 70 ; 80 ].
Os coeficientes de assimetria das quatro distribuições são negativos para o Brasil como um todo : – 0,33 para o Bacharelado, – 0,53 para a Licenciatura, – 0,31 para Sistemas de Informação e – 0,43 para Engenharia de Computação.
As distribuições de todos as habilitações, em todas as regiões, têm o mesmo comportamento, têm cauda mais acentuada no lado esquerdo do gráfico, ou seja, são assimétricas à esquerda.
Na sequência, os resultados verificados para cada uma das questões discursivas de Formação Geral serão apresentados, estabelecendo -se relações com os conteúdos abordados em cada uma delas.
Os comentários da Banca de docentes corretores a respeito do observado na correção das respostas dos estudantes, suas impressões e conclusões serão apresentados junto à análise de cada questão.
Cumpre esclarecer que, tendo em vista que as questões discursivas de Formação Geral são padronizadas, ou seja, constam de todas as provas, os comentários da Banca são os mesmos para todas as carreiras acadêmicas, sendo direcionados a todos os estudantes que participaram do ENADE/2011.
A seguir, serão analisados os desempenhos da Área de Computação ( Bacharelado, Licenciatura, Sistemas de Informação e Engenharia de Computação ) nas duas questões discursivas de Formação Geral do ENADE/2011, comparando os resultados obtidos com comentários para cada questão. 3.3.1.1 Análise da Questão Discursiva 1 do Componente de Formação Geral Os dados de Computação ( Bacharelado, Licenciatura, Sistemas de Informação e Engenharia de Computação ), obtidos a partir das respostas à questão 1, encontram -se na Tabela 3.15 e no Gráfico 3.19.
Nessa questão – de melhor desempenho dentre as duas de Formação Geral – os alunos de todo Brasil tiveram média 57,7.
A maior média para a questão 1 foi obtida na região Norte ( 59,7 ), e a menor, na região Centro-Oeste ( 55,9 ).
Quanto à variabilidade das notas, o desvio padrão de todo o Brasil foi 35,4, o mesmo da região Norte.
O menor desvio padrão foi obtido na região Sul ( 35,0 ) e o maior desvio padrão foi obtido na região Nordeste ( 36,3 ).
As medianas das regiões Sudeste e Centro-Oeste e a do Brasil como um todo foram iguais ( 60,0 ).
Nas regiões Sul e Nordeste a mediana foi 65,0 e na região Norte foi 67,5.
As notas máximas e mínimas da questão discursiva 1 foram as mesmas para todas as regiões do Brasil, respectivamente, 0,0 e 100,0.
O Gráfico 3.19 mostra a distribuição das notas na questão discursiva 1 do componente de Formação Geral para alunos concluintes de Bacharelado, Licenciatura, Sistemas de Informação e Engenharia de Computação.
Observa -se que a maior frequência corresponde aos alunos que obtiveram notas no intervalo ( 90 ; 100 ], para as quatro habilitações.
Também para todos as habilitações a assimetria é negativa nas cinco Grandes Regiões do país.
Tomando -se os alunos do Brasil como um todo, os coeficientes são : – 0,34 para o Bacharelado, – 0,42 para a Licenciatura, – 0,26 para Sistemas de Informação e – 0,42 para Engenharia de Computação.
Nota -se, ainda, que para a habilitação de Licenciatura há uma proporção maior de alunos com notas nos intervalos mais altos. 3.3.1.2 Comentários sobre a correção das respostas à Questão Discursiva 1 De maneira geral, a aplicação da chave de correção da questão 1 de Formação Geral não apresentou qualquer dificuldade digna de menção.
Para isso contribuíram, indubitavelmente, os ajustes feitos após a correção da amostra e a reunião entre todos os membros da Banca de docentes corretores.
As poucas dúvidas, todas pontuais, apresentadas pelos corretores, foram acompanhadas e respondidas pela coordenação e subcoordenação da correção das questões de Formação Geral, por meio da ferramenta de Gerenciamento de Dúvidas do Sistema de Correção On-line.
Não houve registro de qualquer ocorrência que pusesse em xeque o padrão de resposta ou a efetividade e a adequação da chave de correção.
Explica -se : trata -se de questão com comando claro, direto e objetivo ( solicitava -se, basicamente, três vantagens justificadas de cursos a distância ), cujas respostas foram corrigidas por meio da aplicação de um chave de correção testada e aprovada previamente.
Havia absoluta clareza quanto aos critérios de avaliação da correspondência entre as respostas dos estudantes e as possibilidades de vantagens de cursos a distância admitidas como corretas no padrão de resposta oficial, além de gradações explícitas ( e fáceis de aplicar ) dos diferentes níveis de pontuação previstos. respostas dos estudantes e as possibilidades de vantagens de cursos a distância admitidas como corretas no padrão de resposta oficial, além de gradações explícitas ( e fáceis de aplicar ) dos diferentes níveis de pontuação previstos.
Felizmente, portanto, não há reparo a registrar em relação à facilidade de aplicação do padrão de resposta e da chave de correção, e nem em relação à atribuição dos diferentes níveis de pontuação previstos.
Todas as dificuldades que poderiam ter obstado a correta aplicação do padrão de resposta oficial e da respectiva chave de correção foram evitadas por meio dos ajustes feitos após a correção da amostra e fartamente debatidos com toda a Banca.
Digno de nota é que quantidade tão significativa de profissionais envolvidos na mesma tarefa – tanto para a questão 1 quanto para a 2 – tenha apresentado tão poucas dificuldades na execução da correção, em termos proporcionais.
Em suma, a correção da questão 1 da prova de Formação Geral do ENADE 2011 foi exemplarmente bem planejada, servindo -se de padrão de resposta muito bem adequado à questão proposta.
Quanto ao tema desta questão, em particular, a Banca verificou que uma parcela significativa de estudantes evocou experiências bastante concretas e próximas de sua realidade.
Houve várias respostas que indicavam uma vivência pessoal de ensino superior na modalidade Educação a Distância ( EaD ), evidenciando o tom de depoentes nos textos apresentados.
Os estudantes, em sua grande maioria, utilizaram parte considerável do espaço de 15 linhas disponíveis para a resposta – e outra parcela menos significativa dos que não o fizeram demonstraram notável capacidade de atender ao comando da questão de maneira objetiva, curta e, via de regra, correta.
Registre -se, ainda, que foram relativamente poucos os casos de respostas que tenham passado ao largo do tema em pauta na questão.
As capacidades de leitura, de compreensão do comando proposto e de expressão escrita que os estudantes avaliados na edição 2011 do ENADE foram satisfatórias.
Não obstante, seria leviano perder de vista que a qualidade dos textos redigidos em resposta às questões discursivas do Exame ainda está muito aquém do que se espera de concluintes de cursos de ensino superior de todas as regiões do país.
Quanto ao conteúdo das respostas, a Banca constatou boa capacidade, por parte da maioria dos estudantes, de compreensão do tema e do comando da questão.
Foram relativamente poucos os casos de respostas que deixaram de enumerar vantagens da modalidade EaD, e proporcionalmente escassos os estudantes que citaram vantagens não previstas no padrão de resposta.
Foi frequente, a tentativa direta de atender ao comando da questão.
Os erros mais comuns, em relação ao padrão de respostas e à grade de correção, foram fruto do desdobramento em vários “ itens ” daquilo que, de acordo com o padrão de respostas oficial, representava uma única vantagem.
Destacaram -se, neste caso, as respostas que apontavam a flexibilidade de horário e/ou local como duas vantagens distintas daquela modalidade de ensino.
Quanto aos diferentes níveis de pontuação previstos, a maior causa de baixas pontuações foi a ausência de justificativas, e mesmo de argumentação, para uma ou mais das vantagens enumeradas.
Isso demonstra não apenas a objetividade das respostas, coerente com a objetividade do comando da questão ( “ enumere três vantagens de um curso a distância ” ), mas também certa dificuldade de formulação plena de um texto, ou ao menos de parágrafos, em formato dissertativo – mesmo diante de uma média de 5 linhas disponíveis para cada vantagem a enumerar.
Os acertos mais comuns, ou seja, os “ itens ” do padrão de respostas mais frequentemente mencionados foram : ( 1 ) a flexibilidade de horário e/ou local ; ( 2 ) a capilaridade do ensino a distância ; ( 3 ) a democratização do acesso à educação de qualidade ; e ( 4 ) os custos menores que os de cursos presenciais.
Dentre as vantagens previstas no padrão de respostas que foram menos citadas, destacam -se a inclusão de pessoas com comprometimento motor, a qualificação de professores e a troca de experiências entre os participantes.
Foram muito comuns, no entanto, as menções à supostamente maior facilidade de acesso a professores ou/e tutores em cursos superiores a distância. 3.3.1.3 Análise da Questão Discursiva 2 do Componente de Formação Geral A Tabela 3.16 mostra que o desempenho na questão 2 ( média 39,1 ) foi inferior ao obtido na questão de número 1 ( média 57,7 ) para os estudantes de Computação ( Bacharelado, Licenciatura, Sistemas de Informação e Engenharia de Computação, tomados em conjunto ).
A região Norte foi aquela onde a média, nessa questão, foi maior ( 42,7 ), e a de menor média foi a região Centro-Oeste ( 38,2 ).
Quanto à variabilidade das notas, o desvio padrão de todo o Brasil foi 31,5, inferior ao obtido na questão de número 1 ( 35,4 ).
O maior desvio nessa questão foi obtido na região Nordeste ( 32,9 ), enquanto o menor foi obtido na região Sul ( 31,0 ).
As medianas foram 50,0 para o Brasil como um todo e para as regiões Norte, Nordeste e Sudeste.
Para as regiões Sul e Centro-Oeste a mediana foi 45,0.
A nota máxima 100,0 foi obtida em todas as regiões do Brasil, bem como a nota mínima 0,0.
O Gráfico 3.20 mostra a distribuição das notas na questão discursiva 2 do componente de Formação Geral para alunos de Bacharelado, Licenciatura, Sistemas de Informação e Engenharia de Computação.
Observa -se que a maior frequência corresponde aos alunos que deixaram a questão em branco para todas as quatro modalidades de curso.
No Gráfico 3.20 nota -se, também, que são bastante similares os comportamentos das distribuições de notas do Bacharelado, Sistemas de Informação e Engenharia de Computação que possuem coeficientes de assimetria, respectivamente, iguais 0,02, – 0,02 e – 0,05.
Por serem coeficientes muito próximos de zero pode -se considerar que estas distribuições são simétricas, como se observa nos histogramas.
Já para a distribuição de notas dos concluintes de Licenciatura o coeficiente de assimetria é negativo ( – 0,35 ) e nota -se uma cauda um pouco mais prolongada do lado esquerdo do gráfico.
Comparando -se o histograma deste grupo de alunos com os demais vemos que nos intervalos acima da nota 40 eles estão em maior proporção. 3.3.1.4 Comentários sobre a correção das respostas à Questão Discursiva 2 Diferentemente da questão 1, a questão 2, cujo tema aborda políticas públicas para a erradicação do analfabetismo e para o aumento da empregabilidade, permitiu mais amplitude na elaboração das respostas dos estudantes, uma vez que solicita a apresentação de uma proposta.
À falta de objetividade técnica na elaboração da questão – sobretudo no seu comando – correspondeu uma miríade de possibilidades interpretativas.
Registre -se que não raras foram as respostas que se utilizaram de lugares comuns e exortações religiosas/humanitárias/cívicas, contudo, na grande maioria dos casos, essas exortações foram usadas como um encerramento do texto e não comprometeram a resposta.
Muito frequente também foi a confusão entre política pública e política partidária, bem como alguns poucos, confundiam programas educacionais com programas educativos nas TVs.
Assim, destacaram -se como propostas/programas mais recorrentes :  prosseguimento das políticas já existentes, mas com o aumento dos investimentos, normalmente sugerindo a ampliação da bolsa família, da Educação de Jovens e Adultos ( EJA ) e da EaD ;  valorização do magistério, construção de escolas e melhoria das atuais ;  parceria do governo com empresas para manutenção de salas de aula visando a alfabetização de seus funcionários, oferecimento de estágios e redução da carga horária em troca de isenção de impostos ;  parcerias com igrejas e ONGs para criação de espaços de alfabetização ;  escolas itinerantes e alfabetizadores em domicílio, principalmente para pessoas com dificuldade de locomoção, como os idosos e deficientes físicos ;  erradicação do trabalho infantil ;  vinculação da bolsa família não apenas à frequência, mas também e, principalmente, ao resultado obtido pelo aluno na escola ;  revisão das políticas atuais, sendo a mais frequente a extinção da aprovação automática ;  críticas consistentes em relação a modelos didáticos considerados inadequados e desestimulantes para a educação não só de adultos, mas de pessoas de todas as idades.  Destacam -se, ainda, outras sugestões apresentadas :  Creches nas escolas onde os pais estão sendo alfabetizados.  Diminuição de duas horas na jornada de trabalho em empresas para funcionários não alfabetizados, para que possam frequentar a escola.  Campanhas educativas vinculadas aos meios de comunicação.  Aumento do número de escolas noturnas.  Formação específica para professores alfabetizadores.
Quanto à relação entre o analfabetismo e a empregabilidade, deve -se sublinhar que nem todos os estudantes estabeleceram claramente o vínculo entre essas duas situações sociais.
Alguns falaram separadamente de uma e de outra.
Mas a maioria fez referência à necessidade de estudo para “ conseguir um bom emprego com um bom salário ”.
Alguns, em menor número, estabeleceram de forma bastante interessante a questão histórica para a situação do Nordeste ; e, também, a relação entre a pessoa analfabeta/com pouca instrução e a desempregada/empregada em serviços mais pesados e pior remunerados, que não teriam condições de educar seus filhos que, por sua vez, também não teriam melhores oportunidades no mercado de trabalho, identificando a formação de um círculo vicioso e a necessidade de sua interrupção.
Muitos afirmaram que o analfabetismo não é o único responsável pelo desemprego e sim a má distribuição de renda.
Grande também foi o número de alunos que criticaram o resultado da pesquisa, afirmando que a mesma é enganosa, uma vez que considera alfabetizada a pessoa que “ desenha seu nome ”, sendo comum associarem a atual política de alfabetização com ganhos eleitorais.
Foi comum a resposta incompleta, em que o estudante fez apenas a análise das desigualdades/crítica do quadro apresentado, ou só apresentou proposta.
Alguns se limitaram a responsabilizar o governo referindo -se de forma bastante genérica à questão da “ educação ” e “ profissional ”.
Entre os equívocos que mais se repetiram, destacou -se a simples análise dos dados apresentados na tabela que consta do enunciado.
Alguns textos, inclusive, estavam corretos, sem que, no entanto, fosse respondida a pergunta.
Da mesma forma, quando os estudantes partiram para a segunda parte da questão, fizeram referência a vários programas já existentes ou simplesmente disseram que é muito importante que existam projetos para a educação.
Outro equívoco recorrente foi a análise da educação no Brasil como um todo.
Essa análise, apesar de correta, não se referia especificamente à questão do analfabetismo, o que, por conseguinte, levava a sugestões que não eram direcionadas à erradicação do analfabetismo, e sim à melhora da educação no Brasil.
Sendo assim, foram apresentadas propostas como o aumento do número de faculdades, o reforço de alunos do Ensino Médio, a criação de escolas técnicas, etc.
Por outro lado, houve um grande número de redações bem escritas e precisas, no que se refere ao que foi exigido pela questão.
Foram análises equilibradas e sensatas, correlacionando a problemática do analfabetismo com o desemprego e a apresentação de sugestões bastante consistentes. 3.3.2 Componente de Conhecimento Específico Na parte da prova relativa às questões discursivas no componente de Conhecimento Específico ( Tabela 3.17 ), observa -se que a média foi mais baixa do que para as questões discursivas do componente de Formação Geral.
Enquanto no componente de Formação Geral a média para estudantes de Computação de todo o Brasil foi 48,5, na parte de Conhecimento Específico a média foi 3,3.
A maior média deste componente foi obtida pelos estudantes da região Nordeste ( 3,8 ), e a menor, pelos da região Norte ( 2,4 ).
Quanto à variabilidade das notas, o desvio padrão de todo o Brasil foi 9,5.
O maior desvio padrão foi encontrado na região Sul ( 10,1 ), e o menor, na região Norte ( 7,8 ).
A nota máxima 100,0 foi obtida por pelo menos um aluno da região Sudeste, enquanto a menor nota máxima foi encontrada na região Centro-Oeste ( 65,0 ).
Além disso, a nota mínima ( 0,0 ) foi obtida por alunos de todas as regiões do Brasil sem exceção.
A mediana do Brasil e por Grande Região foi 0,0, o que indica que mais da metade dos alunos obteve nota zero, dentre estes estão os que deixaram as questões discursivas do componente de Conhecimento Específico em branco.
O Gráfico 3.21 representa a distribuição das notas nas questões discursivas no componente de Conhecimento Específico segundo a opção de Bacharelado, Licenciatura, Sistemas de Informação e Engenharia de Computação.
As modas destas distribuições ocorrem no intervalo [ 0 ; 10 ], com percentual maior do que 80 \% dos alunos de Licenciatura e Sistemas de Informação neste intervalo.
Nota -se que há pouquíssimas notas nos intervalos subsequentes, acima de 5 \% apenas as distribuições de notas de Engenharia de Computação nos intervalos ( 10 ; 20 ] e ( 20 ; 30 ] e para o Bacharelado no intervalo ( 10 ; 20 ].
Este comportamento das distribuições é confirmado pelos coeficientes de assimetria positivos e muito altos para os quatro cursos, em todas as regiões.
No Gráfico 3.21 os coeficientes de assimetria são : 2,83 para o Bacharelado, 6,23 para a Licenciatura, 7,42 para Sistemas de Informação e 3,85 para Engenharia de Computação.
A análise de cada uma destas questões será feita a seguir. 3.3.2.1 Análise da Questão Discursiva 3 do Componente de Conhecimento Específico Na questão 3, cujos resultados aferidos encontram -se descritos na Tabela 3.18, a média dos estudantes de Computação ( Bacharelado, Licenciatura, Sistemas de Informação e Engenharia de Computação ) de todo o Brasil foi 7,6.
A maior média nessa questão foi obtida pelos alunos da região Nordeste ( 9,2 ), enquanto a menor média foi obtida na região Norte ( 6,0 ).
Quanto à variabilidade das notas, o desvio padrão de todo o Brasil foi 20,0.
O maior desvio padrão foi obtido na região Nordeste ( 22,3 ), enquanto o menor foi obtido na região Centro-Oeste ( 18,2 ).
A nota máxima, 100,0 pontos, foi alcançada por pelo menos um aluno de todas as regiões.
A mediana do Brasil como um todo e a nota mínima foram zero, o que se repetiu para todas as Grandes Regiões, sem exceção.
O Gráfico 3.22 mostra a distribuição das notas na questão discursiva 3, do componente de Conhecimento Específico da Área de Computação para alunos do Bacharelado, Licenciatura, Sistemas de Informação e Engenharia de Computação.
Essa distribuição tem moda nas questões em branco, com percentual maior para os alunos de Licenciatura e Sistemas da Informação, próximo a 70 \%.
Adicionando -se este percentual ao de notas zero, para estes dois cursos, a proporção quase supera 80 \%.
Nota -se, ainda, que para o curso de Engenharia de Computação há uma porcentagem em torno de 3 \% de alunos com notas espalhadas em todos os intervalos.
O coeficiente de assimetria é bastante alto para os quatro cursos, os maiores são 5,59 para a Licenciatura e 5,80 para Sistemas de Informação, já os menores são 1,99 para o Bacharelado e 1,64 para Engenharia de Computação.
A assimetria é positiva para todas as distribuições por região.
O maior coeficiente é 8,30, da distribuição das notas dos alunos de Licenciatura da região Sudeste, e o menor, 0,70, do curso de Engenharia de Computação da região Norte. 3.3.2.2 Comentários sobre a correção das respostas à Questão Discursiva 3 Trata -se de uma questão que solicitava ao estudante a elaboração de dois algoritmos, um recursivo e outro iterativo, para retornar o n-ésimo termo da sequência de Fibonacci.
A definição da sequência é fornecida no enunciado da questão, o que em tese mitiga o problema de o assunto “ sequência de Fibonacci ”, eventualmente, não ter sido visto em todos os cursos para os quais essa questão foi elaborada.
O enunciado estava adequado e continha as informações necessárias e suficientes para a elaboração da solução.
O espaço destinado à apresentação da resposta foi justo, mas não se forneceu algum outro espaço que o estudante pudesse usar como “ margem de manobra ” na eventual ocorrência de um erro na escrita de algum dos algoritmos.
O resultado final obtido na correção foi muito ruim, pois se somarmos as respostas em branco com as respostas com as notas abaixo de 4,0 pontos, teremos um alto percentual de insuficiência.
Não era de esperar que os concluintes tivessem tanta 3.3.2.3 Análise da Questão Discursiva 4 do Componente de Conhecimento Específico A Tabela 3.19 contém as informações relativas à questão 4 do conjunto de questões do componente de Conhecimento Específico.
O desempenho dos estudantes de todo o Brasil nesta questão foi ainda inferior ao desempenho na questão de número 3.
A média geral do Brasil foi 1,9, sendo a menor média registrada na região Norte ( 1,0 ) e a maior na região Sul ( 2,3 ).
A nota máxima ( 100,0 ) foi atingida em todas as regiões.
A mediana e a nota mínima foram zero também em todas as regiões do Brasil, sem exceção. dificuldade nessa questão.
Afinal, a construção de algoritmos desse tipo é um exercício frequente ao longo de todos os semestres, em todos os tipos de curso da área.
As respostas refletiram desconhecimento dos conceitos básicos sobre algoritmos.
Esse resultado corrobora um fato conhecido : a dificuldade do \textbf{aprendizado} de algoritmos e programação no mundo inteiro.
Entretanto, pode -se afirmar que os estudantes que responderam à questão, com nota igual ou superior à nota 4,0, deram soluções com formato semelhante aos apresentados na versão final do padrão de respostas.
Aspecto característico dessas soluções foi o uso da sintaxe das linguagens C ou semelhantes ( principalmente Java ).
Entretanto, foram encontradas, também, soluções construídas com o uso da linguagem Python, o que pode indicar o crescimento dessa linguagem nos currículos dos cursos brasileiros da Área de Computação.
Esse já é um tema comum nas listas de discussão da Sociedade Brasileira de Computação ( SBC ).
É também relevante destacar que, dentre aqueles que tentaram resolver a questão, uma porcentagem significativa de alunos apresentou uma solução para versão iterativa que envolve a definição de um vetor para armazenar todos os termos da sequência de Fibonacci, até o n-ésimo, o que não é necessário para produzir o resultado requisitado pela questão.
É possível atribuir o uso dessa solução à falta de prática dos estudantes em programação.
Excluindo -se o fato notório do desinteresse em resolver as questões propostas, há uma indicação clara de que as nossas metodologias e práticas de ensino de algoritmos e programação precisam ser revistas. 3.3.2.3 Análise da Questão Discursiva 4 do Componente de Conhecimento Específico A Tabela 3.19 contém as informações relativas à questão 4 do conjunto de questões do componente de Conhecimento Específico.
O desempenho dos estudantes de todo o Brasil nesta questão foi ainda inferior ao desempenho na questão de número 3.
A média geral do Brasil foi 1,9, sendo a menor média registrada na região Norte ( 1,0 ) e a maior na região Sul ( 2,3 ).
A nota máxima ( 100,0 ) foi atingida em todas as regiões.
A mediana e a nota mínima foram zero também em todas as regiões do Brasil, sem exceção.
O Gráfico 3.23 representa a distribuição de notas na questão discursiva 4, no componente de Conhecimento Específico para alunos dos quatro cursos da Área de Computação ( Bacharelado, Licenciatura, Sistemas de Informação e Engenharia de Computação ).
A maior frequência, para todas as distribuições ocorre na categoria “ em branco ”, chegando a 70 \% para o Bacharelado e a Engenharia de Computação e superando 75 \% para Licenciatura e Sistemas de Informação.
Observa -se que uma maior proporção de estudantes deixou essa questão em branco, comparando -se à questão discursiva 3.
Como 80 \% ou mais dos alunos não resolveu a questão 4 ou tirou nota zero em todos os cursos, observam -se pouquíssimas notas distribuídas nos demais intervalos.
Os coeficientes de assimetria são positivos e muito altos.
Os maiores são de Licenciatura ( 14,00 ) e Sistemas de Informação ( 18,74 ), os menores são relativos aos cursos de Bacharelado ( 4,98 ) e Engenharia de Computação ( 4,49 ).
Este comportamento se repete em todas as regiões, sendo o menor coeficiente de assimetria igual a 4,10, da distribuição de notas de Bacharelado na região Sul, e o maior, 31,95, de Sistemas da Informação na região Nordeste.
Destaca -se, ainda, que para os cursos de Licenciatura e Sistemas da Informação a proporção de alunos que deixou a questão 4 em branco ou tirou nota zero chega a 99,4 \%. 3.3.2.4 Comentários sobre a correção das respostas à Questão Discursiva 4 A questão versava sobre a transformação de uma estrutura de dados em outra.
Ela requeria certo nível de raciocínio ao exigir uma determinada performance ( tempo linear ) do algoritmo fornecido pelo estudante, sem extrapolar a capacidade de um aluno que tenha uma formação adequada em algoritmos e estruturas de dados.
O espaço destinado à resposta foi adequado.
Alguns necessitaram fazer duas colunas no espaço, mas isso foi apenas porque utilizaram linguagens menos concisas.
O desempenho dos estudantes na questão 4 foi muito ruim, como retratado pelo resultado quantitativo.
Não era de se esperar que os concluintes tivessem tanta dificuldade nessa questão.
Afinal, Estrutura de Dados é uma das disciplinas chave dos cursos de Computação e daí, possivelmente, a escolha do tema para ser explorado em questão discursiva.
Exercícios desse tipo são tão frequentes nos cursos da área quanto os de Algoritmos e Programação.
Nada justifica um resultado tão ruim, a não ser o desinteresse pela prova e a já constatada necessidade de revisão nos processos de ensino e \textbf{aprendizado} de Algoritmos e Programação.
Dentro os que buscaram resolver a questão, os erros mais comuns foram a tentativa de usar uma inserção padrão em árvore, que leva tempo log N e faz com que seja violado o desempenho pedido no enunciado ; o esquecimento do aterramento das folhas da árvore e o não controle correto dos limites de avaliação colocados nas chamadas recursivas.
O fraco resultado é um indicativo claro de uma das seguintes possibilidades :  falta de compromisso por parte dos estudantes, que não se motivaram em fazer a prova ;  falta de compreensão dos conceitos básicos de estruturas de dados ;  falta de capacidade de programar sem o auxílio do que encontram por meio de pesquisa nas ferramentas de busca ;  falta de capacidade de montar algoritmos diferentes do padrão mais divulgado na literatura e em livros textos.
Esta questão é extremamente básica, especialmente o item a, que pede a montagem da estrutura de dados que representa uma árvore em qualquer linguagem de programação ( valendo dois pontos ).
Assim, qualquer nota abaixo de dois indica que o estudante desconhece qualquer linguagem de programação ou mesmo de estruturas de dados.
Entre as questões de baixa pontuação, ressaltamos respostas tais como “ para implementar uma árvore eu usaria uma pilha ”, ou “ a estrutura certa para implementação é uma fila ” e até mesmo “ para implementar uma árvore a estrutura de dados é a recursão ”, que ocorreram centenas de vezes e denotam total falta de compreensão sobre o que são estruturas de dados e de como implementá -las.
Este desconhecimento também está patente até mesmo em respostas corretas.
Houve um grande número de estudantes que descreveu corretamente a estrutura de um nó no subitem ( a ), porém se referiu a ele como"lista duplamente encadeada", provavelmente em função da presença de dois ponteiros em cada nó.
Isso pode indicar uma falta de conhecimento teórico, mostrando que o concluinte percebe semelhanças nas formas das várias estruturas, mas não consegue compreender o arcabouço teórico subjacente ao que está sendo apresentado.
Obviamente, isto é amplamente agravado pela falta de conhecimento prático demonstrado pela manipulação errada da estrutura no algoritmo do subitem ( b ). 3.3.2.5 Análise da Questão Discursiva 5 do Componente de Conhecimento Específico A Tabela 3.20 contém as informações relativas à questão 5 do conjunto do componente de Conhecimento Específico.
O desempenho dos estudantes nessa questão ainda pior do que os observados nas questões 3 e 4.
A nota média dos estudantes de todo o Brasil foi 0,3.
Na região Centro-Oeste a média foi 0,0, indicando que poucos alunos obtiveram pontuação nesta questão.
A maior média foi 0,5, registrada na região Sul.
Quanto à variabilidade das notas, o desvio padrão dos alunos do Brasil, como um todo, foi 4,0.
O desvio padrão, menor do que nas demais questões discursivas, mostra que o fraco desempenho foi mais generalizado nesta questão.
A nota máxima ( 100,0 ) foi alcançada nas regiões Sudeste e Sul, enquanto a menor nota máxima foi registrada na região Centro-Oeste ( 50,0 ).
A mediana e a nota mínima foram zero para todas as regiões do Brasil, confirmando grande quantidade de notas zero na questão.
O Gráfico 3.24 representa a distribuição das notas da questão discursiva 5 no componente de Conhecimento Específico para alunos do Bacharelado e da Licenciatura.
Mais uma vez, destaca -se o grande número de estudantes que deixaram a questão em branco, para os quatro cursos da Área de Computação.
Fato interessante, repetindo a questão 3, é que a maior parte de respostas que receberam notas mais elevadas tenha sido desenvolvida na linguagem C.
Isso não tem nenhuma relação com a linguagem em si, mas possivelmente com o fato de que grande parte das IFES, que possuem um corpo discente intelectualmente elitizado, ainda leciona programação nesta linguagem. 3.3.2.5 Análise da Questão Discursiva 5 do Componente de Conhecimento Específico A Tabela 3.20 contém as informações relativas à questão 5 do conjunto do componente de Conhecimento Específico.
O desempenho dos estudantes nessa questão ainda pior do que os observados nas questões 3 e 4.
A nota média dos estudantes de todo o Brasil foi 0,3.
Na região Centro-Oeste a média foi 0,0, indicando que poucos alunos obtiveram pontuação nesta questão.
A maior média foi 0,5, registrada na região Sul.
Quanto à variabilidade das notas, o desvio padrão dos alunos do Brasil, como um todo, foi 4,0.
O desvio padrão, menor do que nas demais questões discursivas, mostra que o fraco desempenho foi mais generalizado nesta questão.
A nota máxima ( 100,0 ) foi alcançada nas regiões Sudeste e Sul, enquanto a menor nota máxima foi registrada na região Centro-Oeste ( 50,0 ).
A mediana e a nota mínima foram zero para todas as regiões do Brasil, confirmando grande quantidade de notas zero na questão.
O Gráfico 3.24 representa a distribuição das notas da questão discursiva 5 no componente de Conhecimento Específico para alunos do Bacharelado e da Licenciatura.
Mais uma vez, destaca -se o grande número de estudantes que deixaram a questão em branco, para os quatro cursos da Área de Computação.
Nota -se que a soma das frequências dos que deixaram a questão em branco com a dos que fizeram a questão e tiraram zero é ainda mais expressiva do que nas questões discursivas anteriores do Componente de Conhecimento Específico, já que \textbf{ultrapassa} 80 \% do total de alunos avaliados em cada curso.
No caso da Licenciatura apenas 1 dos 833 presentes ( 0,1 \% ) resolveu e obteve nota diferente de zero.
No grupo Sistemas de Informação 99,7 \% dos 12604 alunos ficaram com nota zero, porque deixaram a questão em branco ou apresentaram uma solução errada.
Nesse contexto, é possível observar, também, que as pouquíssimas notas diferentes de zero estão mais presentes no histograma do curso de Engenharia de Computação.
Diante de tal situação os coeficientes de assimetria são excessivamente altos : 12,53 para a distribuição do Bacharelado, 26,87 para Licenciatura, 26,63 para Sistemas de Informação e 10,67 para Engenharia de Computação. 3.3.2.6 Comentários sobre a correção das respostas à Questão Discursiva 5 A questão tratava de tópico específico em todos os seus detalhes e exigia memorização do funcionamento ou da diferença de funcionamento de três tipos de memória cache.
O assunto é específico de apenas uma disciplina, não sendo reforçado pelo uso em outras disciplinas.
O assunto provavelmente não é tratado em todos os cursos e é de difícil compreensão.
O espaço destinado à solução foi inadequado, por ser muito pequeno.
Se o estudante soubesse a resposta teria que gastar algum tempo planejando uma boa forma de descrevê -la.
O resultado geral obtido na questão foi muito ruim, o pior de todas as três questões.
A complexidade da pergunta pode ter assustado os estudantes, parecendo não “ valer a pena ” gastar tempo pensando em uma boa resposta discursiva em comparação com o resto da prova.
Ela exigia tempo para apreciação do que foi pedido e dos detalhes que poderiam ser colocados em tempo hábil.
A pergunta não é adequada ao ENADE e o baixo índice de resposta, mesmo apenas qualitativa, reflete esse fato.
A maioria das soluções corrigidas descrevia a função das memórias cache – que é o que foi guardado do curso – sem qualquer ligação com o que foi pedido na questão.
Os que responderam, o fizeram, quase sempre, de forma apenas qualitativa.
Esse fato foi percebido pelos próprios estudantes ; há várias provas com \textbf{observações} sobre a falta de propriedade da proposição dessa questão, principalmente para o curso de Sistemas de Informação. 3.3.3 Considerações Finais da Banca sobre a correção das questões discursivas do Componente de Conhecimento Específico O resultado deixa claro que o preparo em Algoritmos e Programação não está no nível desejado para os nossos formandos.
Houve um grande número de respostas em branco.
Além da falta de conhecimento, isso indica baixo comprometimento na realização da prova.
Este baixo comprometimento também é evidenciado pelo grande número de estudantes que tirou nota inferior a 40,0 em questões básicas de qualquer curso de Computação ( as questões 3 e 4 ).
No total de questões respondidas, um grande número recebeu nota zero.
Entre essas, deve -se destacar que algumas respostas estavam absurdas.
Grande parte dos graduandos parece não compreender o que é um algoritmo ou uma estrutura de dados.
Isso Outro fato relevante é o grande número de respostas"enciclopédicas"- estudantes que responderam com conceitos decorados, mas que, apesar de corretos, não têm qualquer relação com a questão em si.
Esse tipo de resposta apareceu claramente na análise da questão 5, que aborda memória cache.
O que isso pode significar?
Problema não entendido?
Tentativa de receber alguns pontos?
A falta de tentativas de validação dos resultados, observado nas duas primeiras questões, talvez tenha ocorrido por pouco tempo para responder a prova, mas de qualquer forma é outro indício ( porque não houve um caso sequer de verificação ) da falta de preparo dos alunos na área de programação ou, novamente, de falta de comprometimento com o ENADE.
Foi detectada, ainda, incapacidade de fazer programas que sejam aceitáveis para uso continuado, com os conceitos de legibilidade e autodocumentação bem trabalhados, necessários para qualquer aplicativo “ sério ” e que tenha permanência.
Poucos estudantes demonstraram tal preocupação.
Apesar disso não ter sido foco de avaliação nessa prova ( e não ter causado perda de pontos ), trata -se de conceitos que talvez devam ser reforçados no ensino de Algoritmos e Programação de nossas instituições, e cobrados futuramente.
A soma do número de alunos com notas médias altas dá um percentual insignificante de sucesso ( parcial ou total ).
Isso é um indicativo claro de falta de compromisso por parte dos estudantes, que não se motivaram em fazer a prova ; de falta de compreensão dos conceitos básicos de programação.
Em particular, destaca -se a falta de capacidade de programar sem o auxílio de exemplos encontrados nas ferramentas de busca ; ou a falta de capacidade de montar algoritmos diferentes dos padrões mais frequentes na literatura e livros texto.
Mas, certamente, apesar de todo descaso pela prova, fica claro que há um problema a ser resolvido no processo de \textbf{aprendizado} dos nossos estudantes de Computação na disciplina Algoritmos e Programação, alvo de duas das questões discursivas.

% matched lemmas: aprendizado, observação, ultrapassar
\end{textsample}
